\section{Группа}
	\tikzsetfigurename{LinearAlgebra_Group_}

\begin{definition}[Группа]
    Непустое%
    \sidenote{Требование непустоты здесь, как и далее, в определениях полукольца и кольца~--- дискуссионный вопрос.}
    множество $G$ с заданной на нём бинарной операцией $\circ: {G} \times {G} \to {G}$ называется \emph{группой} $(G ,\circ)$, если выполнены следующие аксиомы:
    \begin{enumerate}
        \item ассоциативность: для любых $a, b, c \in G$ выполнено $(a \circ b) \circ c = a \circ (b \circ c)$;
        \item наличие нейтрального элемента $e$: для любого $a \in G$ выполнено $e \circ a = a \circ e = a$;
        \item наличие обратного элемента: для любого $a \in G$ существует $a^{-1} \in G$, такой что $a \circ a^{-1} = a^{-1} \circ a = e$.
    \end{enumerate}
\end{definition}

Иными словами, группа~--- это моноид с дополнительным требованием наличия обратных элементов.

\begin{definition}[Абелева группа]
    Если операция $\circ$ коммутативна, то говорят, что группа \emph{абелева}.
\end{definition}

По аналогии с полугруппой, нотация $x = y \binopFrom{\algebraic{G}} z$ будет использоваться, чтобы явно указать, что операция $\binop$ взята из моноида $\algebraic{G}$.

\begin{example}
    Рассмотрим несколько примеров групп.
    \begin{itemize}
        \item Целые числа $\BbbZ$ с операцией сложения $+$ являются группой.
              Получается дополнением моноида из предыдущего раздела обратными по сложению элементами.
        \item Целые числа $\BbbZ$ без нуля%
              \sidenote{
                  При наличии нуля возникают трудности с нейтральным элементом.
                  Логично считать $1$ нейтральным по умножению, однако $0 \cdot 1 = 0$, а не 1, как того требует определение.}
              с операцией умножения $\cdot$ не являются группой, так как нет обратных по умножению.
              Действительно, возьмём $a = 3$, тогда должен существовать $a^{-1} \in \BbbZ$, такой что $3 \cdot a^{-1} = 1$.
              Видим, что $a^{-1} = 1/3$, но $1/3 \notin \BbbZ$.
        \item Множество обратимых%
              \sidenote{
                  Квадратная матрица $M$ называется обратимой, если существует матрица $N$, называемая обратной, такая что $M \cdot N = N \cdot M = I$, где $I$~--- единичная матрица.
                  К сожалению, не все матрицы являются обратимыми, потому, чтобы сконструировать группу, нам приходится требовать обратимость явно.}
              матриц с операцией матричного умножения задают группу.
    \end{itemize}
\end{example}
